\documentclass[11pt]{article}
\usepackage[margin=1in]{geometry}
\usepackage{booktabs}
\usepackage{amsmath}
\usepackage{longtable}
\usepackage[hidelinks]{hyperref}

\title{Mass-to-Light Ratios of Binary Galaxies:\\
A Reproduction of Schweizer (1987)}
\author{W.\ Clark}
\date{\today}

\newcommand{\ML}{\ensuremath{\mathcal{M}/L}}

\begin{document}
\maketitle

\begin{abstract}
\noindent
We reproduce the binary-galaxy mass-to-light analysis of Schweizer (1987,
ApJS 64, 427; hereafter S87, with Papers I and II at ApJS 64, 411 and 417).
Recovered \ML{} values agree with her Table~4 to a median ratio of 0.989
(mean 1.004, scatter 8.9\%); all 30 entries fall within her quoted 22\% formal
error. This note is written to be read alongside S87 --- it does not restate
her derivations --- but it specifies everything needed to regenerate our
numbers. Section~\ref{sec:method} summarises the method and why it works,
\S\ref{sec:repro} gives the recipe, \S\ref{sec:corr} lists defects we found in
our own code, and \S\ref{sec:dev} lists every remaining departure from S87,
classified as a deliberate \emph{improvement} or a \emph{liberty} taken where
the paper is silent. Section~\ref{sec:honma} extends the reproduction to
Honma (1999): nine of his published \ML{} values are recovered within 68\%
from his own pair list, but not his flat \ML{}--separation result; his parent
sample is rebuilt to within 5\% from present-day catalogues (RC3, NED via TAP,
HyperLEDA, CfA ZCAT-95) and 38--43 of his 57 pairs are recovered, with the
remainder traced to an open ambiguity in the units of his isolation criterion. Code:
\url{https://github.com/ODU-William-Clark/Binary_87}.
\end{abstract}

\section{Method}
\label{sec:method}

\subsection{Why a pair of galaxies can be weighed}

Two galaxies bound to one another orbit their common centre of mass, and the
speed of that orbit is set by the mass enclosing it. If we could watch a pair
around a full orbit we could weigh it directly. We cannot: each pair is seen
once, at one instant, from one direction. Three things are therefore hidden.
We measure the \emph{projected} separation $r_p$ --- the component of the true
three-dimensional separation $r$ perpendicular to the line of sight, so
$r_p \le r$ --- and only the \emph{radial} velocity difference $\Delta v$, the
component of the relative velocity along the line of sight. And we never see
the orbit's eccentricity, because an instantaneous snapshot does not trace a
path.

The way through is statistical. For a Keplerian pair of total mass $M$ at
separation $r$ with eccentricity $e$ and eccentric anomaly $E$ (the angle
parameterising position around the orbit), the vis-viva relation
$v^2 = GM(2/r - 1/a)$ with $r = a(1-e\cos E)$ gives $v^2 = GM(1+e\cos E)/r$.
Hence
\begin{equation}
\Psi \;\equiv\; \frac{v}{\sqrt{M}}
\;=\; \left[\frac{G\,(1+e\cos E)}{r}\right]^{1/2}
\label{eq:psi}
\end{equation}
depends only on the geometry of the orbit, not on the mass.

\textbf{That is the whole trick.} Because $\Psi$ is mass-free, a Monte Carlo
can predict its distribution --- and the distribution of its projection
$\Psi_p$ --- from assumed distributions of separation and eccentricity alone,
with no knowledge of any masses. The observations supply $\Delta v$ and $r_p$.
The only quantity linking the two is the mass, so demanding that the observed
sample look like a random draw from the predicted distribution \emph{measures}
the mass. Eccentricity cannot be measured per system, so a population
distribution $f(e)$ must be posited instead; S87 tries five, chosen to bracket
the plausible range from circular to nearly radial.

\subsection{The statistic}

Throughout, $\Delta v$ denotes the \emph{observed} line-of-sight velocity
difference and $r_p$ the observed sky-plane separation. The projected,
observable counterpart of eq.~(\ref{eq:psi}) is
\begin{equation}
\Psi_p \;\equiv\; \frac{|\Delta v|}{\sqrt{M}} ,
\end{equation}
and the pair $(r_p, \Psi_p)$ is what can be compared with data.

Masses are parameterised as $M_i = y\,[\,q(t_1)L_1 + q(t_2)L_2\,]_i$ for pair
$i$, where $t$ is de Vaucouleurs morphological type (negative for early types:
$t=-5$ is an elliptical, $t=+5$ an Sc), $q(t)$ is an assumed \ML{} shape, and
$y$ the scaling to be fitted. Luminosities are visual isophotal, in solar
units, so \ML{} is $\mathcal{M}/L_V$ in solar units throughout. One then forms
\begin{equation}
u_i \;=\; \Phi\!\left(\Psi_{p,i}^{\rm obs} \,\middle|\, r_{p,i},\ \mathrm{model}\right),
\qquad
\Psi_{p,i}^{\rm obs} = \frac{|\Delta v_i|}{\sqrt{M_i}} ,
\label{eq:u}
\end{equation}
where $\Phi$ is the cumulative distribution of $\Psi_p$ \emph{conditional on}
that pair's separation, convolved with a Gaussian of width
$\sigma_{v,i}/\sqrt{M_i}$ to carry the velocity errors. Conditioning on $r_p$
matters: $\Psi \propto r^{-1/2}$, so $r_p$ is the strongest single predictor of
$\Psi_p$.

Equation~(\ref{eq:u}) is the probability integral transform: if a value is
drawn from a distribution and mapped through that distribution's own CDF, the
result is uniform on $[0,1]$. So if $y$ is correct the $u_i$ are uniform, and
if $y$ is wrong they pile up at one end. Uniformity is tested with a $\chi^2$
over five equal cells, accumulating the bound-pair probabilities $P(B|C_i)$
--- the posterior probability that pair $i$ is physically bound rather than a
chance projection --- instead of integer counts. Scanning $y$ and minimising
$\chi^2$ gives the best fit.

The five eccentricity distributions are $f_1(e)=\delta(0)$, $f_2=2(1-e)$,
$f_3=1$, $f_4=2e$, $f_5=\delta(0.9)$ (S87 eqs.~43a--43e); the three morphology
functions are $q_1(t)=6.47-0.39t$ (Faber \& Gallagher 1979), $q_2(t)=2$ for
$t\le0$ and $1$ otherwise, and $q_3(t)=1$ (eqs.~46a--46c). Note the ordering
constraint these imply: $\langle 1+e\cos E\rangle$ = 1.000, 0.917, 0.835,
0.750, 0.596 for $f_1 \ldots f_5$, and $M \propto \langle 1+e\cos
E\rangle^{-1}$, so \textbf{\ML{} must increase monotonically from $f_1$ to
$f_5$}. We use this as a convergence test in \S\ref{sec:results}.

\section{Data and reproduction recipe}
\label{sec:repro}

\paragraph{Data provenance --- important.} \texttt{Binary\_gal\_87.txt} is a
direct transcription of S87 Table~1 (and Table~2 for $P(B|C)$), not a
recomputation. Its columns are: pair index; $t_1,t_2$; $\theta$ (arcmin);
$r_p$ (kpc); $|\Delta v|$ (km\,s$^{-1}$); $L_1, L_2$ (units of
$10^{10}L_\odot$); a reference mass; $P(B|C)$; and $\sigma_1,\sigma_2$
(km\,s$^{-1}$, with \texttt{...} for missing). The observational reduction
chain --- B1950 to galactic coordinates, the Local Group apex correction, and
S87 eqs.~(10)--(12) for $L$, $r_p$ and $\Delta v_p$ --- is therefore
\emph{not} exercised by this pipeline.
\texttt{Binary\_87\_luminosity\_calculator.py} implements it and can be used to
audit individual entries, but nothing downstream depends on it. Likewise
\texttt{Binary\_87\_probabilites.py} recomputes $P(B|C)$, but the fit reads
S87's published values from the data file; that script is diagnostic only.
The sample is 48 pairs (96 galaxies); S87 \S Vc excludes the five with
$\Delta v > 2500$\,km\,s$^{-1}$ and $P(B|C)=0.15$ (pairs 18, 22, 23, 36, 40),
leaving $N=43$. Our $P>0.5$ cut removes exactly those five.

\paragraph{Environment.} Python 3 with \texttt{numpy}, \texttt{scipy},
\texttt{pandas}, \texttt{matplotlib}, \texttt{astropy}. The intermediate
\texttt{.npz} is $\approx$106\,MB.

\paragraph{Order of execution.}
\begin{enumerate}\itemsep2pt
  \item \texttt{Binary\_87\_r\_distro.py} --- Lucy deconvolution of the
        lognormal $f(r_p) \propto r_p^{-1}\exp[-0.63(\ln r_p - 4.32)^2]$
        (S87 eqs.~38--39) into the spatial $f(r)$. Imported, not run directly.
  \item \texttt{Binary\_87\_psi.py} --- Kepler orbits and $\Psi$ sampling.
        Imported, not run directly.
  \item \texttt{Binary\_87\_Para\_Main\_2.py} --- Monte Carlo ensembles
        (S87 steps M1--M8) $\rightarrow$
        \texttt{simulated\_psi\_data\_test.npz}. ($\approx$6\,s)
  \item \texttt{Binary\_87\_ML\_8.py} --- the $y$ scan $\rightarrow$
        \texttt{ML\_fit\_results.csv}. ($\approx$140\,s)
  \item \texttt{Binary\_87\_Chi\_Square\_3.py} --- best-fit $y$ per model
        $\rightarrow$ \texttt{best\_fit\_y\_results.csv}.
  \item \texttt{Binary\_87\_calibrate\_chi2.py} --- bootstrap calibration,
        confidence intervals, and the closure test (\S\ref{sec:chi2}).
\end{enumerate}

\paragraph{Configuration.} \texttt{n\_samples} $=5\times10^5$;
\texttt{nb} $=5$ cells; \texttt{NBINS\_PSI} $=50$ with the upper edge at
$1.2\max\Psi_p$ (\texttt{PSI\_RANGE\_PCT} unset); \texttt{USE\_R\_CONDITIONING}
$=$ True with \texttt{N\_R\_BINS} $=8$ quantile bins; \texttt{Y\_GRID\_N}
$=4000$ linear points on $[0.1,500]$; $H_0 =
50$\,km\,s$^{-1}$\,Mpc$^{-1}$; $r \in [10,1500]$\,kpc; Lucy \texttt{n\_iter}
$=5$. \texttt{PSI\_RANGE\_PCT} and \texttt{Y\_GRID\_N} are read from
environment variables of the same name.

\section{Corrections to our own implementation}
\label{sec:corr}

These are not departures from S87. They are defects in our own code;
correcting each moved us back toward her method. We record them because they
changed the answer materially.

\begin{enumerate}\itemsep4pt
  \item \textbf{Projection geometry.} S87 step M7 projects $r$ onto the plane
        of the sky and $\Psi$ onto the line of sight using a \emph{single}
        random orientation. We had drawn two independent angles, uniform in
        angle rather than uniform on the sphere. Corrected:
        $\langle r_p/r \rangle = 0.7856$ (exact $\pi/4 = 0.7854$),
        $\langle (r_p/r)^2 \rangle = 0.6669$ (exact $2/3$),
        $\langle (\Psi_p/\Psi)^2 \rangle = 0.3334$ (exact $1/3$); previously
        $0.637$, $0.500$, $0.500$ --- a 50\% error in the second moment.
  \item \textbf{M8 selection function.} S87 eq.~(25) gives the
        \emph{reciprocal} correction $\nu^{-1}(r_p) \propto \mathrm{dex}(-0.0017 r_p)$
        for physical pairs accidentally excluded, so eq.~(26)'s
        $\nu = c\exp(r_p/255\,\mathrm{kpc})$ implies an acceptance
        $\propto \exp(-r_p/255)$. We had used
        $1 - \exp[(r_p-800)/255]$, nearly flat out to an invented hard wall at
        800\,kpc (no such cutoff appears in the paper). Corrected, the kept
        fraction falls from 92\% to 69.5\%, the simulated $r_p$ median moves
        from 74.1 to 59.5\,kpc against an observed 54.5, and the fraction of
        simulated pairs beyond the observed maximum halves from 13\% to 6.8\%.
        This was the single largest contributor to our residual disagreement
        with Table~4.
  \item \textbf{Velocity errors assigned per pair instead of per galaxy.} S87
        step 5 assigns the mean sample error where no multiple spectrograms
        exist --- per galaxy. A single \texttt{try/except} around
        $\sqrt{\sigma_1^2+\sigma_2^2}$ meant that a pair with one measured
        error had it discarded. Of the 43 pairs, 14 have both errors, 18 have
        exactly one, and 11 have neither; the 18 were wrongly set to
        9\,km\,s$^{-1}$, and the 11 to 9 rather than $\sqrt{2}\times9 = 12.7$.
  \item \textbf{Conditional rather than marginal PDF.} S87 step M8 histograms
        the ensemble jointly in $r_p$ and $\Psi_p$ and slices at each pair's
        $r_{p,i}$. We had used the marginal $\Psi_p$ distribution.
  \item \textbf{Lucy deconvolution.} S87 \S Ve states ``a close approximation
        to $f(r_p)$ was achieved after cycling through five iterations''; we
        used ten. Separately, Lucy's normalisation requires
        $\int_0^r K(r_p|r)\,dr_p = 1$, but our $r_p$ grid began at
        $r_{\min}=10$\,kpc, so the integral was 0.00 at $r=r_{\min}$ and 0.44
        at 11.5\,kpc and $f(r)$ decayed without bound as iterations increased.
        With the grid starting at 0, $P(r<25\,\mathrm{kpc})$ converges
        (0.0339, 0.0303, 0.0288, 0.0287 at 5, 10, 20, 50 iterations). An exact
        $r = r_p$ coincidence also produced \texttt{inf} and is now guarded.
  \item \textbf{Sampled $r$ overwritten} by an independent redraw after
        projection, decoupling the stored $r$ from the stored $\Psi$.
  \item \textbf{$y$ grid too coarse.} At 1000 points the step is 0.50, which is
        11\% of $y$ for the $q_1$ models and tied $f_3$ with $f_4$, breaking
        the monotonicity test of \S\ref{sec:method}. At 4000 points they
        separate (4.23 vs 4.73) and all three $q$ columns are monotonic.
  \item \textbf{Bin width in the calibration script.} With $r_p$ conditioning
        each group has its own histogram range, hence its own bin width. Using
        one group's width for all pairs mislocated $\Psi_p^{\rm obs}$ on the
        CDF grid and biased the recovered $y$ by a factor of $\sim$3.
  \item \textbf{Numerical.} Vectorised Newton--Raphson for Kepler's equation
        (residual $8.9\times10^{-16}$) in place of a generic root finder, and
        the $r$ distribution computed once rather than in every worker process.
\end{enumerate}

\section{Departures from S87}
\label{sec:dev}

\subsection{Deliberate improvements}

\begin{longtable}{p{0.26\textwidth} p{0.68\textwidth}}
\toprule
\textbf{Item} & \textbf{What we do, and why} \\
\midrule
\endfirsthead
\toprule
\textbf{Item} & \textbf{What we do, and why} \\
\midrule
\endhead
\bottomrule
\endfoot

Confidence intervals on $y$ &
We report bootstrap intervals. S87 adopts a single 22\% formal error for all
solutions, from a closure test on 10 circular-orbit samples (\S VI[c]). A
$\Delta\chi^2=1$ interval is \emph{not} available here: because $\chi^2$ steps
whenever a $u_i$ crosses a cell boundary, the sublevel set is disconnected ---
up to five disjoint blocks in our runs --- so its extrema are the convex hull
of a set with holes, not an interval. \\[4pt]

Calibration of $\chi^2_\nu$ &
We calibrate the \emph{minimised} $\chi^2$ by parametric bootstrap rather than
comparing it against unity (\S\ref{sec:chi2}). \\[4pt]

Closure test &
Extended from S87's 10 circular-orbit samples to 300 samples for each of the
five eccentricity models, using the fitted conditional PDFs
(\S\ref{sec:chi2}). \\[4pt]

Degenerate polynomial fits &
S87 \S VI(b)iii fits a second-order polynomial over $\chi^2 \le 3\chi^2_{\min}$
and reports its vertex. She describes her curves as ``nearly parabolic''; ours
are not, and the unguarded fit returned a negative \ML{} for $f_1$. We keep her
window but fall back to the raw $\arg\min$ when the fit degenerates, and quote
the raw $\arg\min$ in \S\ref{sec:results}. \\
\end{longtable}

\subsection{Liberties taken where S87 is silent}

\begin{longtable}{p{0.26\textwidth} p{0.68\textwidth}}
\toprule
\textbf{Item} & \textbf{What we do} \\
\midrule
\endfirsthead
\toprule
\textbf{Item} & \textbf{What we do} \\
\midrule
\endhead
\bottomrule
\endfoot

Monte Carlo ensemble size &
$5\times10^5$ per model. With 8 $r_p$ bins, $10^4$ samples leaves only
${\sim}23$ per $\Psi_p$ bin. \\[4pt]

$\Psi_p$ histogram binning &
50 bins with the upper edge at $1.2\max\Psi_p$. This edge drifts with
\texttt{n\_samples}; a percentile alternative is available via
\texttt{PSI\_RANGE\_PCT} and is statistically safer, but did not improve
agreement with Table~4 and is not the shipped default. \\[4pt]

$r_p$ conditioning resolution &
8 quantile bins. S87 states only that the ensemble is ``histogrammed in $r_p$
and $\Psi_p$''; the grid is not given. \\[4pt]

$y$ grid &
4000 linear points on $[0.1,500]$. S87 gives the range (step 3) but not the
sampling. \\[4pt]

Optical-pair field model &
1000 fields of 50--150 galaxies with magnitudes uniform on $[9,18]$, against
S87's $10^6$ fields with $\log N(<m) = 0.6m + c_i$ (eq.~27) and her isolation
criterion S1 (eq.~5). This is a genuine shortfall, but it is inert: the fit
uses S87's published $P(B|C)$, not ours. Retained as a diagnostic only. \\
\end{longtable}

\subsection{Previously misclassified --- now known to be faithful}

Four items appeared in an earlier version of this note as departures. Checking
them against the paper shows they are S87's own prescriptions, and they are
listed here so the record is not misleading.

\begin{itemize}\itemsep3pt
  \item \textbf{$\sigma_v = 9$\,km\,s$^{-1}$ fallback.} Prescribed by step 5;
        Paper I gives $\langle\sigma_v\rangle \approx 9$\,km\,s$^{-1}$. The
        real defect was applying it per pair (\S\ref{sec:corr}, item 3).
  \item \textbf{Uniform galaxy positions} in the optical-pair simulation. S87
        \S Vc(i): ``The coordinates of the galaxies were chosen from a uniform
        distribution. Since clusters were avoided during the actual search,
        this is a sufficiently good approximation.'' Our earlier claim that
        this was a liberty forced by clustering was contradicted by the paper.
  \item \textbf{Independent sampling of $r$ and orbital phase.} Steps M2, M3
        and M5 draw $e$, mean anomaly, and $r$ independently, even though a
        real orbit links them through $r = a(1-e\cos E)$. This is inherited,
        not introduced --- but see \S\ref{sec:open}.
  \item \textbf{Minimum-separation cut.} S87 bias B5 applies no correction for
        pair overlap. The cut in our code
        (\texttt{APPLY\_MIN\_SEP}) is a corrupted form of her isolation
        criterion S1 (eq.~5), not a B5 cut. It is inert at these separations
        and affects nothing.
  \item \textbf{Kepler solver.} Not a departure at all; S87 M4 specifies only
        that Kepler's equation be solved. Moved to \S\ref{sec:corr}.
\end{itemize}

\section{Reading the \texorpdfstring{$\chi^2$}{chi-squared}}
\label{sec:chi2}

The quantity reported is the \emph{minimum} of $\chi^2$ over the $y$ scan. A
minimum is biased low: one scans thousands of trial values and keeps the best,
so even for a perfect model the smallest result sits well below the mean of the
statistic. It must not be compared against the degrees of freedom.

Two Monte Carlo calibrations, both in
\texttt{Binary\_87\_calibrate\_chi2.py}, make this quantitative. First, the
\emph{unminimised} statistic behaves as expected: drawing $u_i$ uniform with the
real $P(B|C)$ weights gives mean 3.897 and median 3.278, against the analytic
$(\mathrm{nb}-1)\sum P_i^2/\sum P_i = 3.886$. The weighting is therefore not the
problem, and $\nu = 5-1 = 4$ is appropriate (following S87, who charges no
degree of freedom for $y$). Second, a parametric bootstrap at the fitted $y$
gives, for $f_4/q_3$, a minimised $\chi^2$ of mean 1.881 and median 1.425 ---
that is, a \emph{perfect} model here scores $\chi^2_\nu \approx 0.47$ (mean) or
$0.36$ (median), not 1.

S87's own Table~4 values run 0.53--1.02, and she selects the model with the
\emph{lowest} $\chi^2_\nu$, not the one nearest unity. Her best model sits at
0.53, essentially exactly the value our calibration says a perfect fit should
produce. Her values were never anomalously low; unity was simply the wrong
yardstick. Binning must not be tuned toward it.

\paragraph{Closure test.} S87 \S VI(c) generated 10 synthetic samples of 43
bound pairs on circular orbits with known \ML{} and recovered
$(\ML)_{\rm out}/(\ML)_{\rm in} = 1.03 \pm 0.07$, whence her 22\% formal error.
Repeating this with 300 samples per eccentricity model, drawn from the fitted
conditional PDFs, gives ratios of 1.005, 1.018, 1.038, 1.065 and 1.057 for
$f_1 \ldots f_5$ --- consistent with hers, and confirming the pipeline is
unbiased. The implied per-sample relative error is $\approx$28\%, somewhat
larger than her 22\%.

\section{Results}
\label{sec:results}

Best-fit $y$ is the raw $\arg\min$ of the scan. Bootstrap intervals are 16th to
84th percentile.

\begin{center}
\begin{tabular}{lrrcrr}
\toprule
Model & $y$ & $\chi^2_\nu$ & 68\% interval & S87 $\mathcal{M}/L$ & ours \\
\midrule
$f_1$ (circular)    & 11.85 & 0.62 & 9.1--14.9  & 13 & 11.9 \\
$f_2 = 2(1-e)$      & 13.85 & 0.68 & 10.7--17.6 & 15 & 13.9 \\
$f_3 = 1$           & 16.48 & 0.31 & 13.0--22.4 & 21 & 16.5 \\
$f_4 = 2e$          & 21.10 & 0.37 & 16.7--29.2 & 24 & 21.1 \\
$f_5 = \delta(0.9)$ & 28.35 & 0.37 & 22.2--39.3 & 32 & 28.4 \\
\bottomrule
\end{tabular}
\end{center}

\noindent
(The last two columns are the $q_3$ model, for which $y$ \emph{is} \ML{}.)
\textbf{Every one of S87's $q_3$ values falls inside our 68\% bootstrap
interval.} Across all 30 entries of her Table~4 --- the three $q$ columns
evaluated at her tabulated morphological types --- the median ratio of ours to
hers is 0.989, the mean 1.004, and the scatter 8.9\%; all 30 lie within her
22\% formal error. For her preferred model $f_4$ under $q_1$ at $t=+2.6$ we
obtain 25.8 against her 25.

Monotonicity in eccentricity (\S\ref{sec:method}) now holds in all three $q$
columns: $q_1$ gives 2.73, 3.48, 4.23, 4.73, 6.73; $q_2$ gives 11.85, 13.85,
16.48, 24.23, 30.85; $q_3$ as tabulated. Earlier runs violated it in two
columns, which was proof the $\arg\min$ had not been located.

Following S87 we take $f_4(e) = 2e$ as preferred, giving
$\mathcal{M}/L_V \approx 21$ for an Sc galaxy and $\approx 40$ for an
elliptical at $H_0 = 50$, in solar units. Since $M \propto H_0^{-1}$ and
$L \propto H_0^{-2}$, $\ML \propto H_0$: multiply by ${\sim}1.4$ for
$H_0 = 70$. We do \emph{not} claim to discriminate between eccentricity models:
the observed $\chi^2_{\min}$ sits at the 48th to 78th percentile of its
bootstrap null for every model, so $\chi^2$ separates none of them.

\section{Extension: reproducing Honma (1999)}
\label{sec:honma}

Honma (1999, ApJ 516, 693; hereafter H99) re-did the binary-galaxy \ML{}
measurement twelve years after S87 with a larger, deliberately wider-separation
sample (57 pairs, $R_p$ out to 400 kpc) and a maximum-likelihood fit in place
of the binned $\chi^2$. His headline conclusion is not the \ML{} value itself
but its \emph{constancy with separation}: ``$\mathcal{M}/L$ does not show any
tendency of increase, but [is] found to be almost independent of the
separation of pairs beyond 100 kpc,'' whence ``the typical halo size of spiral
galaxies is less than $\sim$100 kpc.'' We reproduced him in two stages:
(i) his \ML{} analysis from his own published pair list, and (ii) his sample
selection from present-day catalogues. Stage (i) reproduces every one of his
twelve published \ML{} entries that can be computed from public data, but
\emph{not} his separation-dependence result (\S\ref{sec:h99results}); stage
(ii) reproduces his parent sample and a core of his pairs, and the remainder
is an open ambiguity in his isolation criterion (\S\ref{sec:h99sel}). This
section was revised after three independent audits against the paper; the
audit findings are incorporated below rather than listed separately. Code is
in \texttt{Binary\_99/} (mirrored under \texttt{honma1999/} in the repository).
Everything is at $H_0 = 50$, as in H99.

\subsection{His data}
\label{sec:h99data}

H99 Table~1 is VizieR catalogue J/ApJ/516/693 (\texttt{honma\_table1\_full.dat},
114 rows = 57 pairs, one row per galaxy): name, partner, galactic coordinates,
heliocentric velocity and error, morphological type, $B$ magnitude, and the
pair-level luminosity-corrected $R_p$, $V_p$ and $(\ML)_p$.

\paragraph{Unit convention (verified).} All separations and velocities are
normalised by $L_{10}^{1/3}$, $L_{10} = (L_1+L_2)/10^{10}L_\odot$ (his
eqs.~6--7), so $R_p = r_p/L_{10}^{1/3}$, $V_p = |\Delta v|/L_{10}^{1/3}$, and
$(\ML)_p = R_p V_p^2/(G\times10^{10})$ with
$G = 4.30091\times10^{-6}$\,kpc\,(km\,s$^{-1}$)$^2\,M_\odot^{-1}$; this
reproduces his $(\ML)_p$ column.

\paragraph{Velocity frame (verified).} H99 says ``heliocentric velocities''
and applies no Local Group correction anywhere. His tabulated velocities match
RC3 heliocentric values to MAD 13\,km\,s$^{-1}$, against 133 for the Yahil et
al.\ correction.

\paragraph{Parse validation.} \texttt{Honma99\_table1\_to\_pairs.py} recomputes
$R_p$ and $V_p$ from the raw columns (flux-weighted mean velocity as the
``luminosity centre,'' $D = \bar v/H_0$, $M_{B,\odot} = 5.48$). Ratios to his
columns: $R_p$ median 1.001, $V_p$ median 1.001. His \S4.1 subgroup split
(27/12/18 pairs at $R_p <100$, 100--200, $>200$\,kpc) falls out of the
recomputed $R_p$ exactly. His photometry is the RC3 corrected system:
$B_{\rm H99} - \mathtt{btc}$ (HyperLEDA) has median $+0.09$, MAD 0.13 over 113
members, against $-0.48$, MAD 0.32 for raw \texttt{bt}.

\subsection{His \texorpdfstring{\ML}{M/L} method}
\label{sec:h99method}

\textbf{Bound pairs} are point masses whose separation vector follows a
\emph{volume} density $\nu(R) \propto R^{-\gamma}$, $\gamma = 2.6$,
$R \ge R_{\min} = 10$\,kpc ($\nu$ enters his Jeans equation, eq.~8, in the
Binney \& Tremaine sense; as a check, projecting the scalar
$p(R) \propto R^{2-\gamma}$ matches his observed $R_p$ at KS $p = 0.96$, with
the KS statistic peaking at his fitted $\gamma = 2.6$, while the scalar
reading is excluded at $p \sim 10^{-29}$). His projection kernel, eq.~12,
carries a spurious $2/\pi$; we use the normalised isotropic kernel. Velocities
from the Jeans equation, $G' = G\times10^{10} = 43009.1$ in luminosity-corrected
units:
\begin{align}
\beta = 0:\quad & \sigma_r^2 = \frac{G'(\ML)}{(\gamma+1)R}, &
\beta = -\infty:\quad & |V|^2 = \frac{G'(\ML)}{R},\ V\perp R, &
\beta = 1:\quad & \sigma_r^2 = \frac{G'(\ML)}{(\gamma-1)R}.
\end{align}
H99 states only the dispersions, not a distribution function, and describes
the ensemble as ``one million pairs \ldots assuming random orientation of
orbital planes.'' We draw an isotropic 3-D Gaussian at $\sigma_r$ and
\textbf{reject speeds above escape}, $v^2 > 2G'(\ML)/R$ --- scale-free, since
$v^2/v_{\rm esc}^2 = \chi^2_3/2(\gamma+1)$ is independent of \ML{} --- so the
ensemble is simulated once at $\ML = 1$ and rescaled. This bound DF is the
physically consistent choice for a bound-pair ensemble; the unbounded
Gaussian used in a first version runs 10--16\% low on every isotropic case
(\S\ref{sec:h99results}). Circular orbits are projected with one shared line
of sight for $R$ and $V$.

\textbf{Optical pairs:} $p_{\rm opt} \propto [1+\xi(r)]\,R_p$ on the window
$R_p, V_p \in [0,400]$ (uniform in $V_p$ for $q=0$; falling with $V_p$ for
$q=1.8$), $\xi = (r_0/r)^q$, $r_0 = 10$\,Mpc,
$r^2 = r_p^2 + (\Delta v/H_0)^2$ in physical units. Whether $r$ uses per-pair
$L_{10}$ or $L_{10}=1$ is immaterial ($\xi \gg 1$ across the window; results
identical).

\textbf{Likelihood} over $(\ML, f)$, $f$ the true-pair fraction (his
eqs.~17--19):
\begin{equation}
\log\mathcal{L} = \sum_i \int g_i(V)\,
\log\!\left[f\,p_{\rm bin}(R_{p,i}, V) + (1-f)\,p_{\rm opt}(R_{p,i},V)\right]dV ,
\label{eq:h99like}
\end{equation}
with $g_i$ a Gaussian in $V_p$ of width $\sigma_{V,i} =
\sqrt{e_1^2+e_2^2}/L_{10}^{1/3}$, normalised to unit weight per pair. This
\emph{expected} log-likelihood --- the observed pair spread as a cloud and
integrated against $\log p$ --- is the literal reading of eq.~17 ($\sum n\log p$
over the phase space) with eq.~18 ($n = \sum_i g_i$). It is not the usual
convolved form $\sum_i \log\int g_i\,p\,dV$; the decisive evidence for the
expected-log reading is the circular-orbit case, whose model has a hard cutoff
at $V_c$ so that part of each error cloud sits where $\log p \to -\infty$: the
convolved form gives 12.0 against his 28, the expected-log form 27.5. (A
first version of this note justified the reading by matching the median
$(\ML)_p$; that check happens to work for the isotropic case and fails for
the circular one, and has been withdrawn.) Eq.~19 as printed has a
$(2\pi\sigma_i)^{-1/2}$ prefactor, which under eq.~18's collective
normalisation would weight pair $i$ by $\sqrt{\sigma_i}$; that reading raises
the isotropic sample-I values to 38.6 and 46.8 when combined with the bound
DF, overshooting his 35 and 36, and is treated as a typographical
$\sigma_i^2$.

\textbf{Sample III} (``30 pure spiral pairs \ldots two spiral galaxies later
than Sa'') is reproduced by the rule ``both members anything but E, S0/SB0,
Pec,'' i.e.\ Sa \emph{and} the irregulars (Im, IBm, Sm) count. Three
independent confirmations: it is the only one of 64 exclusion combinations
giving $N=30$; it gives a spiral fraction of 70.2\% for sample I against his
stated 70\%; and it gives a mean raw separation of 205.8\,kpc for sample III
against his stated $\sim$206\,kpc (\S4.1).

Implementation liberties, all tested insensitive: $R_{\max} = 2000$\,kpc
(800--5000 identical); $4\times10^6$ pairs vs his $10^6$; model density per
pair from a slab $|R_p - R_{p,i}| < 5$\,kpc (5--50 identical) on a
2\,km\,s$^{-1}$ grid smoothed at 3\,km\,s$^{-1}$ (1--3 identical); errors
floored at 1\,km\,s$^{-1}$ (0.1--5 identical; only $\sigma=0$, one pair,
matters); grids of 90 geometric points on $[2,150]$ and 71 in $f$ (grid
spacing $\sim$5.7\%, so quoted values carry $\pm3\%$ resolution).

\subsection{Results against his tables}
\label{sec:h99results}

Bound DF, unit weights. Intervals: 1-D profile at $\Delta\log\mathcal{L}=0.5$,
and the 2-D-contour extent at $\Delta\log\mathcal{L}=1.15$, which is the
convention of his likelihood figures.

\begin{center}
\begin{tabular}{llrrcccl}
\toprule
sample & orbits & $q$ & \ML & $f$ & 68\% (1-D) & 68\% (2-D) & H99 \\
\midrule
I (57)  & isotropic & 0   & 31.8 & 0.88 & 26.2--40.5 & 22.6--46.8 & $35^{+7}_{-5}$, 0.88 \\
I       & isotropic & 1.8 & 35.0 & 0.77 & 28.8--46.8 & 24.9--54.2 & $36^{+8}_{-4}$, 0.71 \\
I       & circular  & 0   & 27.5 & 0.87 & 24.9--30.3 & 23.7--33.3 & $28^{+5}_{-3}$, 0.88 \\
I       & circular  & 1.8 & 28.8 & 0.73 & 26.2--31.8 & 24.9--35.0 & $30^{+6}_{-3}$, 0.73 \\
III (30) & isotropic & 0  & 14.6 & 0.96 & 11.5--18.6 & 9.9--22.6  & $15^{+5}_{-3}$, 0.95 \\
III     & isotropic & 1.8 & 14.6 & 0.90 & 11.5--19.6 & 9.9--23.7  & $16^{+4}_{-4}$, 0.84 \\
III     & circular  & 0   & 10.9 & 0.95 & 9.4--12.6  & 9.0--14.6  & $12^{+4}_{-3}$, 0.95 \\
III     & circular  & 1.8 & 10.4 & 0.87 & 9.4--12.6  & 8.6--14.6  & $12^{+3}_{-3}$, 0.86 \\
I       & radial $\beta=1$ & 0 & 49.2 & 0.89 & 38.6--65.8 & 33.3--79.8 & $42^{+34}_{-7}$ \\
\bottomrule
\end{tabular}
\end{center}

All nine published values lie inside the 2-D 68\% intervals, and all but one
inside the 1-D intervals; circular orbits agree to 2--9\%, isotropic to
3--9\%. The four Sample~II entries (109 pairs at $a=1.5$) cannot be
reproduced because he published only the sample-I pair list. Note the $f$
values are within his quoted errors but not always close (0.77 vs 0.71, 0.90
vs 0.84 for the clustered cases).

\paragraph{His separation-dependence result is not reproduced.} H99 \S4.1
splits sample~I at $R_p = 100$ and 200\,kpc (27/12/18 pairs) and finds
$\ML = 36^{+10}_{-5}$, $37^{+31}_{-17}$, $25^{+29}_{-13}$ --- flat, which is
his evidence that halos are smaller than the pair separations. On the same
split (counts reproduce exactly) our fit gives \textbf{23.7} (19.6--31.8),
36.7 (23.7--56.8), 38.6 (19.6--76.1): a \emph{rising} trend, with the inner
bin below his 68\% interval and $f$ pinned at the 0.99 boundary there.
(Monte Carlo realisations move these by $\pm1$ grid step, $\sim$6\%.) The
discrepancy is confined to $R_p < 100$\,kpc, survives every variant tested
(bound vs Gaussian DF, slab width 5--50\,kpc, the $\sqrt{\sigma_i}$
weighting), and is therefore not a resolution or DF artefact. It is the one
result of his we cannot reproduce, and it bears directly on his physical
conclusion; we flag it rather than explain it.

\paragraph{Free cross-checks that pass.} Seven of his 57 pairs lie above the
$(\ML)_p = 20$ envelope of his Fig.~1b, matching his ``expected number of
optical pairs \ldots about 7'' and our $f = 0.88$ ($\approx 6.8$ optical
pairs).

\subsection{Reconstructing his sample: data sources and access}
\label{sec:h99sources}

H99 built his parent from ``NED, supplied with RC3'': 6,475 galaxies with
$B \le 15.5$, $1000 < cz < 4500$\,km\,s$^{-1}$, $|b| > 20^\circ$ (he writes
``$b \le 20^\circ$ excluded''; his Table~1 has pairs at $b = -36^\circ$, so
$|b|$ is meant). Four sources approximate that today:

\begin{itemize}\itemsep3pt
  \item \textbf{RC3} (local copy): positions, heliocentric velocities,
        $B_T^0$ and error. 23,010 rows; 10,662 velocities; only 4,014
        magnitude errors --- requiring the error (as an earlier finder did)
        cuts the parent to 1,357.
  \item \textbf{NED via TAP} (\texttt{ned\_tap\_fetch.py}):
        \url{https://ned.ipac.caltech.edu/tap}, table \texttt{NEDTAP.objdir}
        (\texttt{prefname, gallon, gallat, prefphytype, z, z\_bibcode}).
        Synchronous queries need \texttt{REQUEST=doQuery} and die at 60\,s;
        asynchronous jobs are \emph{also} aborted at $\sim$60\,s, so the
        velocity shell ($cz = 840$--4676, 43,572 galaxies) was fetched as
        eight $\Delta z = 0.0016$ slices with pauses: nine requests. TAP has
        no photometry; NED's legacy per-object CGI (wrapped by
        \texttt{astroquery}) currently fails server-side (``EGRET error'').
        Page scraping must not be used.
  \item \textbf{HyperLEDA} (\texttt{atlas.obs-hp.fr/hyperleda/fG.cgi},
        table \texttt{meandata}; header space-separated, rows tab-separated):
        every galaxy with \texttt{btc} $\le 15.8$ in 12 longitude slices ---
        102,800 objects, 88,281 with a velocity. \texttt{btc} is the RC3
        corrected-$B$ system H99 used (\S\ref{sec:h99data}), available for all
        of them, which is what makes a magnitude-thresholded companion pool
        possible. It also supplied $B$ for the 13 members of his pairs that
        RC3 lacks (NGC 2979 has no $B$ anywhere modern; his 13.6 is used,
        flagged).
  \item \textbf{CfA ZCAT, June 1995} (VizieR VII/193; \texttt{zcat95\_parse.py},
        B1950$\to$galactic via \texttt{astropy} FK4): 41,600 galaxies with
        velocities, a genuine pre-1999 compilation.
\end{itemize}

\paragraph{Defining the 1999 epoch.} A galaxy is redshift-known in 1999 if it
positionally matches (within $0.05^\circ$; ZCAT positions are coarse) an RC3
row with a native velocity, a ZCAT-95 entry, or a NED galaxy whose preferred
redshift cites a $\le 1999$ bibcode; 28,596 of 102,800 qualify. Others are
redshift-unknown: ineligible as members, blocking by projection only. This is
a lower bound on the 1999 catalogue (NED supersedes preferred redshifts; 90\%
of current ones in the shell cite post-1999 sources, including three pairs in
his own Table~1).

\subsection{His selection criteria and how we read them}
\label{sec:h99sel}

\texttt{Honma99\_finder\_leda.py EPOCH MAGCOL FRAME COMPMAG ISO A} implements
H99 \S2.1. For each criterion: the text, the reading adopted, the evidence,
and its status.

\begin{enumerate}\itemsep4pt
  \item \textbf{Members:} \texttt{btc} $\le 15.0$, $|b| > 20^\circ$,
        $1000 < v_{\rm helio} < 4500$. Like-for-like parent at his 15.5
        definition (which includes companion-only galaxies): \textbf{6,821
        vs his 6,475}; members at 15.0: 6,139. (An earlier version compared
        6,137 members with his 6,475 --- different populations.) His two
        $\sim$7\% data-quality cuts are pair-level, after the blind step
        (``if redshift uncertainty is not reported for any of the two galaxies
        of pair \ldots''); requiring HyperLEDA \texttt{e\_v} and \texttt{e\_bt}
        of members removes 2 galaxies and no pairs, so his cuts are
        \emph{not} reproduced --- HyperLEDA's error availability is not
        NED-1999's. \emph{Status: faithful except the quality cuts, which are
        unreproducible.}
  \item \textbf{Pair cuts are luminosity-corrected:} $m_{1+2} \le 13.5$,
        $R_p \le 400$, $V_p \le 400$ with the $L_{10}^{1/3}$ scaling.
        13 of his pairs exceed 400\,kpc raw and 3 exceed 400\,km\,s$^{-1}$
        raw. $L_{10}$ from \texttt{btc}, $M_{B,\odot}=5.48$, $D=\bar v/H_0$.
        No $(1+z)$ factor (he works in km\,s$^{-1}$ throughout; an earlier
        version applied one, worth 0.8\%). \emph{Status: settled.}
  \item \textbf{Isolation logic.} Printed: a pair is kept ``if all of its
        companion galaxies brighter than $m_{1+2}+2.0$ mag satisfy both'' 5 and
        6. Taken strictly (every companion far in projection \emph{and} far in
        velocity; harmful if close in either) this keeps 0 of his 57 pairs,
        since criterion~6 alone has no spatial bound. The operational reading
        --- a companion is harmful iff inside the search cylinder, close in
        both projection and velocity; redshift-unknown companions, close in
        projection --- follows from his next sentence, ``Parameters $a$ and
        $b$ determine the volume for companion search,'' and from his \S3.2
        example of a projected foreground/background galaxy. Known-redshift
        galaxies outside the shell are not sample members and are ignored
        (treating them as blind blockers keeps 38 of 57 rather than 43).
        \emph{Status: settled, but by the ``volume'' sentence, not by the
        printed criterion.}
  \item \textbf{Isolation units: OPEN.} As typeset, both thresholds scale
        with $L_{10}^{1/3}$, and H99 states the design explicitly: ``the
        volume depends only on the total luminosity of a pair.'' Scored on
        his own 57 pairs, the raw reading (1000\,kpc, 600\,km\,s$^{-1}$)
        keeps 43--45 against 23--26 for the literal one --- but this test
        cannot discriminate: every one of his pairs has $L_{10} \ge 1.30$, so
        the scaled thresholds are uniformly stricter on his sample and any
        looser rule keeps more of it. Diagnostics that are \emph{not} monotone
        in permissiveness split:
        \begin{center}
        \begin{tabular}{lccc}
        \toprule
        & literal scaled & raw & H99 \\
        \midrule
        final pairs, $a=2.5$ & 78 & 328 & 57 \\
        blind-step rejection & 43\% & 23\% & $\sim$30\% \\
        pairs dropped as groups & 38 & 171 & ``two groups'' \\
        of his 57 recovered & 18 & 38 & --- \\
        Sample II ($a=1.5$) pairs & 255 & 512 & 109 \\
        II/I ratio & 3.27 & 1.56 & 1.91 \\
        \bottomrule
        \end{tabular}
        \end{center}
        Sample size, group count and blind rate favour the literal reading;
        recovery and the II/I ratio favour raw. Either reading requires his
        1999 catalogue to have differed from ours, in opposite directions.
        His \S3.2 ``8 Mpc'' remark corresponds to 400, not $b\times400 = 600$
        km\,s$^{-1}$; $b = 1$ changes the scores by at most two pairs.
        \emph{Status: open; the finder runs both.}
  \item \textbf{Companion magnitudes.} His only statement on magnitudes
        (``corrections \ldots according to RC3'') and the photometry of his
        own members point to the corrected system, so \texttt{btc} is the
        textually supported choice and the funnel-faithful one
        (blind rate 23\%). The blockers of five of his pairs sit 0.01--0.09
        mag above the \texttt{btc} threshold, and three of those five
        (NGC 5916A, NGC 5169, UGC 11027) fall below it on raw \texttt{bt} ---
        so his blind search plausibly returned raw magnitudes for some
        objects. Raw \texttt{bt} everywhere recovers 40 but drops the blind
        rate to 7\%; a hybrid (\texttt{btc} for RC3 galaxies, raw otherwise)
        recovers 43 at 9\%. \emph{Status: inference, bracketed; \texttt{btc}
        is the default.}
  \item \textbf{Redshift-blind rejection:} criterion 5 only, on primary
        candidates; his ``about 30\%'' vs our 23\% (\texttt{btc}).
        \emph{Faithful.}
  \item \textbf{Group rejection:} any galaxy in two or more surviving pairs
        marks a group; its pairs are dropped. Wording matches; scale does not
        (171 vs his ``two possible groups'' --- a symptom of item 4), and the
        rule removes two of his own pairs. \emph{Faithful in form.}
  \item \textbf{Not his rule:} mutual nearest neighbours (89\% of his pairs
        vs 54\% of extras; partial only).
\end{enumerate}

\subsection{Reconstruction results}

\begin{center}
\begin{tabular}{lllrrrr}
\toprule
mode & isolation & companions & parent$^\dagger$ & pairs & of his 57 & extras \\
\midrule
1999, helio & raw    & \texttt{btc} & 6,821 & 328 & 38 & 290 \\
1999, helio & raw    & hybrid       & 6,821 & 447 & 43 & 404 \\
1999, helio & scaled & \texttt{btc} & 6,821 &  78 & 18 &  60 \\
2026, all $v$ & raw  & \texttt{btc} & 7,480$^\ddagger$ & 481 & 40 & 441 \\
\bottomrule
\end{tabular}
\end{center}
\noindent {\footnotesize $^\dagger$ like-for-like at $B\le15.5$; $^\ddagger$ members at 15.0.}

Of his 57 pairs, 56 have both members in HyperLEDA and 55 pass every parent
cut; all losses occur at isolation (15) and group rejection (2). Where a pair
is matched, our $R_p$ agrees with his to a few per cent.

\paragraph{What the misses and extras are.} Under \texttt{btc} companions,
five misses have a bright in-shell companion inside 1000\,kpc and
600\,km\,s$^{-1}$: two (NGC 7185/7188 with NGC 7180 at 170\,kpc; NGC
7444/7443 with NGC 7450) block under either magnitude system and are
violations of his stated criterion however read; three block only under
corrected magnitudes (item 5). Two misses are member photometry differences
(ESO 122-IG002: \texttt{btc} 15.10 vs his 14.5; NGC 2979). One blocker is an
\texttt{SDSSJ} object absent in 1999. The extras are systematically wider and
faster than his pairs (median $R_p$ 172 vs 94\,kpc, $V_p$ 47 vs 26). His
sample spans $\bar v = 1513$--3653\,km\,s$^{-1}$ with none beyond 3653 despite
a 4500 cut; H99 gives the reason himself --- ``the number of bright pairs
which satisfy the criteria 3 and 4 becomes small at large redshift'' --- and
28\% of our 1999-epoch members and 94 of our 328 pairs sit beyond 3653, so
this is a flux-limit effect his procedure handled more aggressively than ours,
not a redshift-completeness effect. (An earlier version of this note
attributed it to 1999 catalogue depth; that is withdrawn.) The extras are thus
tied to the open isolation question, not to a missed criterion.

\subsection{Catalogue drift}

Present-day NED both adds members (post-1999 redshifts) and removes them (new
galaxies de-isolate old pairs). With NED shell galaxies as blockers and no
magnitude threshold, 2026 blockers would kill 289 of 293 extras and 39 of his
own 42 recovered pairs; 1999-epoch blockers 183 and 21. A modern study is a
fresh selection on a version-stamped catalogue, not an extension of his
sample.

\subsection{Recipe}

Scripts in \texttt{Binary\_99/}; \texttt{numpy, scipy, pandas, astropy}.
\begin{enumerate}\itemsep2pt
  \item \texttt{ned\_tap\_fetch.py} per $z$ slice ($\times8$) $\to$
        \texttt{ned\_shell\_galaxies.csv}.
  \item HyperLEDA 12-slice pull (README) $\to$ \texttt{leda\_btc15.8\_allsky.csv}.
  \item VII/193 \texttt{zcat.dat}; \texttt{zcat95\_parse.py} $\to$
        \texttt{zcat95\_galactic.csv}.
  \item \texttt{leda\_photometry\_fill.py} $\to$ \texttt{rc3\_photometry\_supplement.csv}.
  \item \texttt{Honma99\_table1\_to\_pairs.py} $\to$ \texttt{honma99\_pairs.csv}.
  \item \texttt{Honma99\_ML.py [bound|gauss] [unit|sqrtsig]} ($\sim$5\,min)
        $\to$ the table of \S\ref{sec:h99results} and the separation bins.
  \item \texttt{Honma99\_finder\_leda.py 1999 btc helio btc raw 2.5} (and
        \texttt{scaled}, and \texttt{1.5}); \texttt{Honma99\_crossmatch.py}.
  \item \texttt{Honma99\_isolation\_test.py}, \texttt{Honma99\_mnn\_test.py},
        \texttt{Honma99\_miss\_audit.py}.
\end{enumerate}

\subsection{Departures from H99, classified}

\begin{longtable}{p{0.30\textwidth} p{0.50\textwidth} p{0.14\textwidth}}
\toprule
\textbf{Item} & \textbf{What we do} & \textbf{Class} \\
\midrule
\endfirsthead
\toprule
\textbf{Item} & \textbf{What we do} & \textbf{Class} \\
\midrule
\endhead
\bottomrule
\endfoot
Heliocentric frame; L-corrected pair cuts; volume-density $\nu$; expected-log
likelihood; Sa and irregulars in sample III; harmful-iff-close-in-both
isolation logic & Readings fixed by his text or his tables. & Settled reading \\[4pt]
Isolation thresholds raw vs $L^{1/3}$-scaled & Diagnostics split; his text
favours scaled; both are run. & Open \\[4pt]
Bound-truncated isotropic DF & He states no DF; bound is the consistent
choice and moves every isotropic case toward him. & Liberty (motivated) \\[4pt]
Companion magnitude system & His text supports corrected; three of his pairs
need raw; bracketed. & Liberty (bracketed) \\[4pt]
Unit-weight error clouds & Eq.~19's $(2\pi\sigma_i)^{-1/2}$ read as a typo;
the literal weighting overshoots. & Settled reading (typo) \\[4pt]
Isolation scored against his pairs; funnel by blocker type; epoch bracketing;
positional cross-match; separation-bin refit; radial case & Machinery he did
not need. & Improvement \\[4pt]
Bulk catalogue access (NED TAP slices, HyperLEDA fG.cgi, ZCAT-95) &
Replaces interactive NED; no scraping. & Improvement \\[4pt]
$R_{\max}$, ensemble size, slab and smoothing widths, grids, error floor,
$0.05^\circ$ tolerance, $M_{B,\odot}=5.48$ & Unstated by H99; tested
insensitive. & Liberty \\[4pt]
His pair-level 7\% quality cuts & Not reproducible from HyperLEDA's error
availability; omitted. & Omission \\
\end{longtable}

\section{Open issues}
\label{sec:open}

\begin{enumerate}\itemsep4pt
  \item \textbf{The $\chi^2(y)$ curve is jagged.} $\chi^2$ steps whenever a
        $u_i$ crosses one of the five cell boundaries, giving many local minima
        within $\Delta\chi^2 < 2$ and a disconnected sublevel set. Percentile
        histogram ranges, Savitzky--Golay smoothing and coarser $y$ grids all
        failed to stabilise it; a finer grid plus the corrected selection
        machinery did, but fragility remains. The principled fix is a statistic
        continuous in $y$ --- Cram\'er--von Mises or Anderson--Darling on the
        $u_i$, or a direct likelihood
        $\sum_i \log f(\Psi_{p,i} \,|\, r_{p,i}, y)$.
  \item \textbf{Independent $(r, E)$ sampling is unquantified.} It is S87's own
        prescription, but a real orbit links $r$ and phase, and conditioning on
        an observed $r$ should skew the phase distribution toward apocentre.
        Because the fit conditions on $r_p$, any resulting tilt does not average
        out. The effect has been identified but not measured; it is the largest
        unexamined physical assumption in the analysis.
  \item \textbf{The deconvolved $f(r)$ is not smooth.} Local logarithmic slopes
        computed from it are noisy at the 10\,000-bin resolution used, which
        limits how far it can be interrogated and would need addressing before
        quantifying issue 2.
  \item \textbf{The observational reduction chain is unaudited.} Our inputs are
        transcribed from S87 Table~1 rather than recomputed
        (\S\ref{sec:repro}), so eqs.~(10)--(12), the apex correction and the
        arcmin/arcsec handling do not enter our results and have not been
        independently checked.
\end{enumerate}

\end{document}
